\documentclass[12pt]{article}

\usepackage{amsmath,amssymb,amsthm}
\usepackage{enumerate}
\usepackage{geometry}
\usepackage{newtxtext}
\usepackage{newtxmath}

\newtheorem{definition}{Definition}
\newtheorem{theorem}{Theorem}
\newtheorem{lemma}{Lemma}
\newtheorem{example}{Example}
\newtheorem{corollary}{Corollary}

\geometry{margin=1in}

\title{Interpolation Theory 1}
\author{-}
\date{\today}

\begin{document}

\maketitle


\begin{theorem}
Let $\lambda \in (0,1)$. Suppose $f : [0,1] \to \mathbb{R}$ is continuous and
satisfies the following condition: for every $0 \le a < b \le 1$, the value
of $f$ at the weighted point $(1-\lambda)a + \lambda b$ is equal to either
the minimum or the maximum of $f$ on $[a,b]$. Then $f$ is constant.
\end{theorem}

\begin{proof}
Suppose for contradiction that $f$ is not constant. Since $f$ is continuous
on the compact interval $[0,1]$, it attains a global maximum value $M$ and a
global minimum value $m$ (there is a video for this, I love this video), and our assumption forces $M > m$ since we want to derive a contradiction.

Define the sets
\[
  A = f^{-1}(M) = \{x \in [0,1] : f(x) = M\},
  \qquad
  B = f^{-1}(m) = \{x \in [0,1] : f(x) = m\}.
\]
Both sets are nonempty (because $f$ attains its extrema) and closed (because
$f$ is continuous, and preimages of closed sets are closed). Pick any point $a_0 \in A$ and any point $b_0 \in B$;
without loss of generality assume $a_0 < b_0$ again for the sake of contradiction.

Build the two sequences inductively. At every stage we maintain points
$a_n \in A$ and $b_n \in B$ with $a_n < b_n$ based on our hypothesis. Given such a pair, consider
the interval $[a_n, b_n]$. Because $a_n \in A$ and $b_n \in B$, the maximum
of $f$ on $[a_n, b_n]$ is exactly $M$ and the minimum is exactly $m$.
Applying the hypothesis to this interval, the weighted interior point is given by the convex combination
\[
  p_n = (1-\lambda)\,a_n + \lambda\,b_n
\]
must satisfy $f(p_n) \in \{m, M\}$. We update the pair according to which
extremum is hit, this feels awfully similar to a binary search:
\begin{itemize}
  \item If $f(p_n) = M$, then $p_n \in A$; we set $a_{n+1} = p_n$ and
        $b_{n+1} = b_n$, so the new interval is $[p_n, b_n]$, which has
        length $(1-\lambda)(b_n - a_n)$.
  \item If $f(p_n) = m$, then $p_n \in B$; we set $a_{n+1} = a_n$ and
        $b_{n+1} = p_n$, so the new interval is $[a_n, p_n]$, which has
        length $\lambda(b_n - a_n)$.
\end{itemize}
In either case the length shrinks by a factor of at most
$\rho := \max(\lambda, 1-\lambda)$, which is strictly less than one because
$\lambda \in (0,1)$ and it cannot be $0$ or $1$. After $n$ steps the length satisfies
\[
  b_n - a_n \;\le\; \rho^n (b_0 - a_0) \;\longrightarrow\; 0.
\]
Therefore $a_n$ and $b_n$ converge to the same limit, call it $x^*$.
Since $A$ is closed and every $a_n$ belongs to $A$, we have $x^* \in A$,
so $f(x^*) = M$. Since $B$ is closed and every $b_n$ belongs to $B$, we
also have $x^* \in B$, so $f(x^*) = m$. This forces $M = m$, contradicting
our assumption that $M > m$. Hence $f$ must be constant.
\end{proof}

\begin{corollary}
If $f : [0,1] \to \mathbb{R}$ is continuous and satisfies
\[
  f\!\left(\frac{a+b}{2}\right) \in
  \left\{\min_{x \in [a,b]} f(x),\; \max_{x \in [a,b]} f(x)\right\}
  \quad \text{for every } 0 \le a < b \le 1,
\]
then $f$ is constant.
\end{corollary}

\begin{proof}
This is the special case $\lambda = \tfrac{1}{2}$ of the theorem, since
$(1 - \tfrac{1}{2})\,a + \tfrac{1}{2}\,b = \tfrac{a+b}{2}$ is the midpoint
of $[a,b]$.
\end{proof}

\end{document}